\section{A ramification cost hidden in small Selmer lengths}
\label{sec:ramification-budget}

The simultaneous twisting theorem supplies a promising ingredient for other
arithmetic problems. Its numerical output requires particular care: the
lengths are normalized by $\lambda_{\mathcal O}(\mathcal O/\ell)=1$,
while the coefficient ring $\mathcal O$ is allowed to become more ramified.
We prove an exact constraint connecting that ramification to the error
budget. It identifies a further obstruction to using the twisting theorem
as a direct proof of Selmer vanishing.

Throughout this section, $\ell$ is a prime, $K/\mathbf Q_\ell$ is a finite
extension, $\mathcal O$ is its ring of integers, $e=e(K/\mathbf Q_\ell)$,
and $v$ is the valuation normalized by $v(\ell)=1$. Thus the positive
values on $\mathcal O$ are positive integer multiples of $1/e$.
For a finite-length $\mathcal O$-module $M$, put
\[
 \lambda_{\mathcal O}(M)
   =\frac{\operatorname{length}_{\mathcal O}M}{e}.
\]
This is exactly the normalization used by the source, rather than the
integer-valued composition length.

\begin{proposition}[The exact discreteness threshold]
\label{prop:length-threshold}
If $M$ is a finite-length $\mathcal O$-module, then
\[
 \lambda_{\mathcal O}(M)<\frac1e\quad\Longrightarrow\quad M=0.
\]
The strict inequality cannot be replaced by a weak inequality. Moreover,
there are nonzero finite modules over coefficient rings of increasing
ramification whose normalized lengths tend to zero.
\end{proposition}

\begin{proof}
A nonzero finite-length module has integer composition length at least one.
For a uniformizer $\varpi$, the nonzero module
$\mathcal O/(\varpi)$ has length one and normalized length $1/e$.
For the final assertion take
\[
 K_h=\mathbf Q_\ell(\zeta_{\ell^h}),\qquad
 \mathcal O_h=\mathbf Z_\ell[\zeta_{\ell^h}],\qquad
 M_h=\mathcal O_h/(\zeta_{\ell^h}-1),\qquad h\geq1,
\]
with $\zeta_{\ell^h}$ primitive. The polynomial
$\Phi_{\ell^h}(1+X)$ is Eisenstein at $\ell$: its constant term is $\ell$,
and modulo $\ell$ it is $X^{\ell^{h-1}(\ell-1)}$.
Consequently $e_h=\ell^{h-1}(\ell-1)$ and $\zeta_{\ell^h}-1$ is a
uniformizer. Therefore
\[
  M_h\ne0,\qquad
  \lambda_{\mathcal O_h}(M_h)=\frac{1}{\ell^{h-1}(\ell-1)}
       \longrightarrow0.
\]
\end{proof}

\begin{lemma}[Valuations of local character values]
\label{lem:root-valuation}
If $\zeta\in\overline{\mathbf Q}_\ell$ has exact order $\ell^h$ with
$h\geq1$, then
\[
 v(\zeta-1)=\frac1{\ell^{h-1}(\ell-1)}.
\]
For $0\leq k<h$ one has
\[
 v(\zeta^{\ell^k}-1)=\ell^k v(\zeta-1).
\]
If $\zeta\in\mathcal O$, the right hand side is at least $\ell^k/e$.
\end{lemma}

\begin{proof}
In residue characteristic $\ell$, an $\ell$-power root of unity reduces
to one. If $a$ is prime to $\ell$, the quotient
$(1-\zeta^a)/(1-\zeta)$ reduces to $a$ and is a unit. Thus all factors in
\[
 \Phi_{\ell^h}(1)
 =\prod_{a\in(\mathbf Z/\ell^h\mathbf Z)^\times}(1-\zeta^a)=\ell
\]
have the same valuation. There are $\ell^{h-1}(\ell-1)$ factors,
proving the first formula. Apply it also to $\zeta^{\ell^k}$, whose
exact order is $\ell^{h-k}$, to obtain the second. Finally
$\zeta-1$ is a nonzero element of the maximal ideal of $\mathcal O$,
so $v(\zeta-1)\geq1/e$.
\end{proof}

\begin{theorem}[Ramification forced by a simultaneous twisting budget]
\label{thm:ramification-budget}
Let $Q$ be a nonempty finite set of finite places of a CM field, disjoint
from its complex conjugate $\overline Q$. Let $\psi$ be a finite-order
character of $\ell$-power order, with all its values in the same
coefficient ring $\mathcal O$. Fix a tame inertia generator $\gamma_w$
at each place of $Q\cup\overline Q$. Suppose, for an integer $m\geq1$,
that:
\begin{enumerate}
\item every $\psi^j$, $1\leq j\leq m$, is nontrivial on tame inertia
      at every place of $Q\cup\overline Q$;
\item for each such $j$,
\[
 \sum_{w\in Q\cup\overline Q}
       v\bigl(\psi^j(\gamma_w)-1\bigr)\leq\varepsilon.
\]
\end{enumerate}
If $k$ is the unique nonnegative integer satisfying
$\ell^k\leq m<\ell^{k+1}$, then
\begin{equation}\label{eq:epsilon-e-budget}
 \boxed{\quad \varepsilon e\ \geq\ 2|Q|\ell^k.\quad}
\end{equation}
\end{theorem}

\begin{proof}
Use the permitted exponent $j=\ell^k$. For every $w$, write
$\zeta_w=\psi(\gamma_w)$. Nontriviality of $\psi^{\ell^k}$ on tame
inertia says that $\zeta_w$ has order $\ell^{h_w}$ for some $h_w>k$.
Lemma~\ref{lem:root-valuation} gives
\[
 v(\zeta_w^{\ell^k}-1)
  =\ell^k v(\zeta_w-1)\geq\frac{\ell^k}{e}.
\]
There are $2|Q|$ places because conjugation is a bijection and the two
sets are disjoint. Summing proves
$\varepsilon\geq2|Q|\ell^k/e$.
\end{proof}

The source theorem is anticyclotomic. Its ramification clause is written
only at $Q$, but at conjugate places the inertia characters are inverses
up to conjugacy. Their orders, and their valuations after subtracting
one, agree. Thus the first hypothesis above follows from the source's
ramification clause for each power. Alternatively, the finite sum in
the second hypothesis already forbids a trivial tame value, since
$v(0)=+\infty$. These observations require the auxiliary places to be
away from $\ell$, as in the source, so the character's ramification is
tame.

\paragraph{What this does, and does not, rule out.}
Suppose the source's estimate for a particular twisted Selmer module is
\[
 \lambda_{\mathcal O}\bigl(S(T\otimes\psi)\bigr)\leq\beta+\varepsilon.
\]
Even in the favorable case $\beta=0$, inequality
\eqref{eq:epsilon-e-budget} implies
\[
 \beta+\varepsilon\geq\frac{2|Q|\ell^k}{e}\geq\frac2e.
\]
The upper bound can therefore never be made smaller than $1/e$ while
the same nonempty ramification set and local smallness conditions are
retained. Integer rounding of this bound cannot certify Selmer
vanishing. This is a limitation of the numerical certificate, not a
proof that the actual Selmer module is nonzero. It neither contradicts
the source nor excludes stronger bounds obtained using additional
information about the Selmer module. The elementary modules in
Proposition~\ref{prop:length-threshold} show why numerical smallness
alone is logically insufficient.

For instance, if $\ell=3$, $m=10$, and $|Q|=2$, then $k=2$ and
$\varepsilon e\geq36$. A budget $\varepsilon=10^{-3}$ forces
$e\geq36000$. The formal resemblance of $10^{-3}$ to a number smaller
than one gives no vanishing conclusion: the minimum positive
normalized length is then at most $1/36000$.

\begin{proposition}[A sufficient additional descent theorem]
\label{prop:fixed-module-descent}
Let $\mathcal O_0$ have absolute ramification index $e_0$, and let
$S$ be a fixed finite-length $\mathcal O_0$-module. Suppose that for
arbitrarily small $\eta>0$ there exist a finite extension of coefficient
rings $\mathcal O_0\subset\mathcal O_\eta$, a finite-length
$\mathcal O_\eta$-module $N_\eta$, and an $\mathcal O_\eta$-linear injection
\[
 S\otimes_{\mathcal O_0}\mathcal O_\eta\ \lhook\joinrel\longrightarrow\ N_\eta
 \quad\text{with}\quad
 \lambda_{\mathcal O_\eta}(N_\eta)\leq\eta.
\]
Then $S=0$.
\end{proposition}

\begin{proof}
If the relative ramification index is $r_\eta$, extension of scalars
multiplies composition lengths by $r_\eta$, while absolute
ramification indices satisfy $e_\eta=r_\eta e_0$. By the structure
theorem for finite modules over a DVR,
\[
 \lambda_{\mathcal O_\eta}
   (S\otimes_{\mathcal O_0}\mathcal O_\eta)
 =\lambda_{\mathcal O_0}(S).
\]
The injection implies $\lambda_{\mathcal O_0}(S)\leq\eta$ for
arbitrarily small $\eta$, so this fixed length is zero.
\end{proof}

This gives a concrete arithmetic research target: an injective
comparison from one fixed ordinary Selmer group into the favorable
twisted groups, with all local conditions matched. The repository's
twisting theorem does not assert such an injection. Changing the twist
changes the representation and its Selmer problem. One cannot use
the invariance of normalized length under coefficient extension as
though twisting were itself coefficient extension.

\subsection{A bounded-support refinement available from the construction}
\label{sec:bounded-twist-support}

There is also a positive extraction. Under exactly the data and
hypotheses of the source's general twisting theorem, its proof gives
an integer $B\geq1$, depending on the fixed Selmer data, Stickelberger
systems, initial coefficient ring and torsion level, but independent
of $m$ and $\varepsilon$, such that the chosen characters may satisfy
\[
 1\leq |Q|\leq B
\]
in addition to all its stated conclusions for every
$m\geq1$ and every $\varepsilon>0$.

Here is the precise proof extraction. The theorem
\texttt{exists\_twistPolynomials} first chooses $s_0$, independently of
$m$ and $\varepsilon$, such that all its output data have
$2|Q|\leq s_0$. Its construction takes $Q=J$, where $J$ is a set of
tied primes containing a designated first prime $P_0$. Consequently
$Q$ is nonempty. The roots-of-unity step then chooses $r,h_0,N$ and
the character values; it does not add auxiliary primes. Thus one may
take $B=s_0$, or $\lfloor s_0/2\rfloor$ with the inherited stronger
cardinality information. To retain nonemptiness in a Lean statement,
one should add a nonempty-set field to the intermediate
\texttt{TwistPolynomials} structure (or return that property alongside
the structure) and carry the proof \texttt{first\_mem} through the two
construction lemmas. This is a source-level proof extraction; we
have not compiled this strengthened interface.

Since every one of $\psi,\ldots,\psi^m$ is ramified at a place of the
nonempty $Q$, the order of $\psi$ exceeds $m$. Taking unbounded $m$
therefore produces infinitely many distinct characters with bounded
\emph{cardinality} of ramification support and simultaneous finite
Selmer-length bounds. Their orders and conductors are not bounded;
the actual auxiliary primes can vary. This is a useful reusable
ingredient beyond the final zero-location theorem. It is not a new
claim of a positive-density twist family, a fixed-conductor result,
or a solution of Goldfeld's conjecture.

\paragraph{Precise source anchors.}
All paths below are relative to \path{grh/src/} at the pinned upstream commit.
\begin{itemize}
\item \path{Lib/General/NumberTheory/PadicLimitsValuation.lean},
      lines 302--338: normalized length and its valuation formula.
\item \path{New/Congruences/Selmer/Line.lean}, lines 130--144:
      \texttt{TwistedLengthLe} quantifies over all sufficiently large
      coefficient rings, which must contain the character values.
\item \path{New/Congruences/Selmer/GeneralTheorem.lean},
      lines 55--76 and 392--396: $\ell$-power order, ramification,
      the normalized Selmer bound, the sum over $Q\cup\overline Q$,
      and all powers from one through $m$.
\item \path{New/Congruences/Selmer/Line.lean},
      lines 646--681; and
      \path{New/Congruences/Selmer/GeneralTheorem.lean},
      lines 403--416: the cardinality bound is chosen before the
      power range and error-dependent choices.
\item \path{New/Congruences/Selmer/Line.lean}, lines 610--616:
      the constructed auxiliary set is $J$; and
      \path{New/Congruences/Selmer/AuxiliaryPrimes.lean},
      lines 429--437: the tied-prime data include $P_0\in J$.
\end{itemize}

The propositions and Theorem~\ref{thm:ramification-budget} are proved
here as elementary consequences of local valuation theory. They are
new additions to this analysis, with no claim that they are new
theorems in the literature. The bounded-support assertion is an
extraction of information already in the repository's construction.
